% =====================================================================
\chapter{Conclusion et perspectives}
\label{ch:conclusion}
% =====================================================================

\section{Rappel de la problématique}

Ce stage partait d'un constat simple : les données publiques de risque
climatique existent, mais la jointure entre l'exposition d'un bâtiment et sa
vulnérabilité n'existe nulle part. Un souscripteur qui doit évaluer un dossier
le fait aujourd'hui à l'œil, en croisant deux portails, sans traçabilité.

La question posée était donc : peut-on construire, à partir de sources
publiques uniquement, un service qui réponde à \og ce bâtiment est-il exposé, et
qu'est-ce qui le rend vulnérable ? \fg{}, avec une qualité d'information
suffisante pour qu'un professionnel de l'assurance s'appuie dessus ?

\section{Réponse apportée}

La réponse est oui, avec des conditions. Le stage a produit un service
fonctionnel, déployé et vérifié, qui :

\begin{enumerate}[leftmargin=1.6em]
  \item résout une adresse française et interroge les sources en parallèle ;
  \item teste le point résolu contre les périmètres d'aléas par géométrie
        réelle, pas par approximation communale ;
  \item associe à chaque aléa la description verbatim du bâtiment issue de la
        BDNB ;
  \item porte, sur chaque fait, sa source, sa date et son niveau de résolution ;
  \item refuse explicitement de produire un score ou une recommandation ;
  \item distingue trois états — présent, absent, inconnu — de sorte qu'une
        panne ne se lise jamais comme une absence de risque ;
  \item met le résultat en situation sur une carte, avec un moteur de scénarios
        fondé sur les classes réglementaires, sans jamais inventer un niveau ;
  \item est couvert par 357 tests automatisés hors-ligne.
\end{enumerate}

Les conditions sont aussi importantes que la réponse. Le service n'est pas
commercialisable en l'état : l'authentification est simulée, le magasin de lots
est en mémoire, la certification attend sa qualification humaine, et le budget
de latence n'est pas encore calibré. Le prototype démontre une thèse ; il ne
remplace pas un produit.

\section{Les trois idées à retenir}

\subsection{Le refus du score est une position tenable}

L'intuition initiale — produire un score global de risque — a été délibérément
abandonnée. Ce n'est pas une concession, c'est un positionnement : sous
Solvabilité~II, un score tiers engage son producteur, se fait reconstruire ou
refuser par l'utilisateur, et mélange des aléas qui ne se compensent pas.

Le produit a été reconstruit sur d'autres fondations — provenance et résolution
— et s'est révélé démontrable sans son chiffre. Le champ de résolution, qui
n'était qu'un attribut parmi d'autres, est devenu le signal de qualité central.

\subsection{L'honnêteté sur l'incertitude est une fonctionnalité}

La trichotomie \code{true}/\code{false}/\code{null}, la distinction entre
absence et panne, l'étiquetage du mode cinématique : tout cela relève du même
principe. Un produit de données ne se juge pas seulement sur ce qu'il affirme,
mais sur la manière dont il dit \og je ne sais pas \fg.

C'est aussi ce qui rend le produit utilisable en contexte réglementaire : une
décision peut être justifiée si l'on sait quel fait a été établi à quel niveau.

\subsection{La soustraction est un travail d'ingénierie}

La phase la plus productive du stage a été celle où le périmètre a été réduit.
Supprimer quatre connecteurs, une couche de notation, un tableau de bord et des
milliers de lignes a augmenté la valeur du livrable. La règle \og pas de code
dormant \fg{} et l'historique Git comme unique archive ont rendu cette
soustraction possible sans perte d'information.

\section{Perspectives}

Le travail identifié pour la suite se répartit en trois horizons.

\subsection{Court terme — finir et qualifier}

\begin{itemize}[leftmargin=1.4em]
  \item \textbf{Qualifier la certification.} Confronter les fixtures des
        quatre aléas certifiés aux cartes officielles et signer la version~1.
  \item \textbf{Calibrer la latence.} Exécuter la campagne de mesures sur au
        moins 100 appels réels et décider, sur données, si la parallélisation
        des couches WFS est nécessaire.
  \item \textbf{Exécuter le test manuel de bout en bout.} Diagnostiquer une
        adresse certifiée TRI un jour de pluie et valider la boucle complète.
  \item \textbf{Terminer le mode VFX} — ou acter qu'il reste un démonstrateur
        partiel, ce qui est défendable puisque la partie gouvernance est livrée.
\end{itemize}

\subsection{Moyen terme — rendre le service exploitable}

\begin{itemize}[leftmargin=1.4em]
  \item \textbf{Authentification réelle} — remplacer le mode démonstration par
        un système de comptes, condition de tout usage externe.
  \item \textbf{Persistance du store de lot} — remplacer le magasin en mémoire
        par une file de traitement, seule façon de tenir le volume.
  \item \textbf{Cache partagé} — sortir le cache du processus pour qu'il
        survive au serverless.
  \item \textbf{Source vectorielle pour le RGA} — câbler une source
        polygonale pour le retrait-gonflement des argiles, ce qui exige un
        amendement constitutionnel, et supprimerait la dernière étiquette
        \code{commune-level-estimate} structurelle.
\end{itemize}

\subsection{Long terme — le produit complet}

\begin{itemize}[leftmargin=1.4em]
  \item \textbf{La trajectoire climatique.} C'est l'argument de fond du
        produit : la non-stationnarité du risque. Le prototype en pose la
        structure — horizons explicites, scénarios étiquetés, unités réelles —
        mais il reste à industrialiser la lecture des projections et à exposer
        les variables par aléa et par horizon dans le contrat.
  \item \textbf{La livraison en lot pour l'actuaire.} Le second acheteur
        attend de la donnée structurée à l'échelle d'un portefeuille, avec les
        mêmes sémantiques que le produit interactif. L'invariant \og le lot est
        N fois l'unitaire \fg{} a été posé précisément pour rendre cette
        extension non coûteuse.
  \item \textbf{La certification étendue.} Passer de quatre aléas certifiés à
        l'ensemble des aléas couverts par un périmètre vectoriel.
\end{itemize}

\section{Mot de la fin}

\typhoon{} est un prototype, et il faut le présenter comme tel. Mais c'est un
prototype qui démontre quelque chose de précis : \emph{la donnée publique
française permet de construire une information de risque climatique par
bâtiment, traçable et défendable} — à condition d'accepter de ne pas tout
savoir et de le dire.

Le stage m'a laissé une conviction : dans un domaine où la donnée est
abondante, la valeur ne se crée pas en ajoutant des sources ou des indicateurs,
mais en établissant avec rigueur une jointure que personne n'avait faite, et en
assumant explicitement les limites de ce que l'on affirme.

\begin{quote}
\itshape
Un produit de données honnête ne dit pas seulement ce qu'il sait.\\
Il dit aussi comment il le sait, et à quel point il le sait.
\end{quote}
